%==============================================================================
%  CAPÍTULO XI
%==============================================================================
\chapter{El incidente de mala fe del deudor concursal}
\label{cap:mala-fe}

\fichaCapitulo{Límites éticos a la extinción de saldos insolutos.}{Causales tasadas del
artículo 169 A, sustanciación procesal, valoración bajo sana crítica y recursos.}

%==============================================================================
\bloqueTeorico
%==============================================================================

\subsection{Insolvencia fortuita e insolvencia oportunista}

El \discharge{} descansa sobre una premisa moral que la \Lreforma{} \citeyearpar{ley-21563}
hizo explícita: no toda insolvencia merece el mismo trato. La crisis del deudor honesto que
fracasó pese a su diligencia ---la insolvencia fortuita--- justifica plenamente el
\freshstart{}: es el riesgo ordinario de la actividad económica, y la sociedad gana
reinsertando a ese deudor antes que condenándolo a la marginalidad. La crisis del deudor que
se sirvió de la insolvencia ---que ocultó activos, distrajo bienes o mintió en sus
declaraciones--- es de otra naturaleza: extinguir sus deudas premiaría el fraude y erosionaría
la confianza sobre la que descansa el crédito.

El incidente de mala fe es el instrumento que separa ambas situaciones. No niega el
\discharge{} como regla: lo condiciona a que la conducta del deudor haya sido honesta. Es,
en palabras simples, el filtro ético del perdón concursal.

\begin{foco}[El equilibrio que busca la figura]
Un sistema sin incidente de mala fe extinguiría por igual la deuda del deudor honesto y la del
oportunista, destruyendo el incentivo a la probidad. Un sistema con un incidente demasiado
laxo ---que denegara el \discharge{} ante cualquier error--- vaciaría el \freshstart{} de
contenido. El arte de aplicar el \art{169 A} consiste en no confundir la omisión inocente con
la distracción dolosa.
\end{foco}

\subsection{Un control civil de probidad patrimonial}

El incidente del \art{169 A} es una \destacar{sanción civil}, no penal. Opera dentro del
propio procedimiento concursal, se tramita como incidente, se prueba conforme a la sana
crítica y su consecuencia es la denegación de un beneficio ---la extinción de saldos---, no una
pena.

Esto lo distingue nítidamente de los delitos concursales del \cpen{}, catalogados hoy por la
\Ldeleco{} \citeyearpar{ley-21595} (capítulo \ref{cap:laboral-penal}). Ambos pueden coexistir
sobre los mismos hechos, pero son independientes: el incidente puede prosperar sin condena
penal, y su estándar ---la sana crítica--- es menos exigente que el de la certeza penal.

\begin{alerta}[Dos consecuencias sobre los mismos hechos]
Un mismo acto de ocultamiento puede fundar, a la vez, el incidente civil de mala fe y la
persecución penal por delito concursal. No hay \emph{non bis in idem} entre ambos, porque
sancionan cosas distintas ---la denegación de un beneficio civil frente a la pena---. El deudor
debe ser advertido de que enfrentar el incidente no lo pone a salvo del proceso penal, ni a la
inversa.
\end{alerta}

%==============================================================================
\bloqueNormativo
%==============================================================================

\subsection{Las hipótesis tasadas del artículo 169 A}

El \art{169 A}, incorporado por la \Lreforma{} \citeyearpar{ley-21563}, tasa las causales
de mala fe en los siguientes términos:

\begin{enumerate}
  \item Falsedad o carácter incompleto de los antecedentes documentales presentados por
  el deudor conforme a los \art{115} o \art{273 A} para abrir la liquidación.
  \item Destrucción u ocultamiento de información patrimonial e instrumentos de comercio
  en los dos años previos al concurso o durante su tramitación.
  \item Actos de distracción u ocultamiento de bienes dentro de los dos años anteriores
  al concurso o durante el mismo.
  \item Sentencia firme que acoja una acción revocatoria especial en contra del deudor.
  \item Condena firme por alguno de los delitos concursales del \cpen{}.
\end{enumerate}

El catálogo precedente corresponde literalmente al texto vigente. El numeral 5 remite a
los delitos concursales del Párrafo 7 del Título IX del Libro Segundo del \cpen{}.

\subsection{Oportunidad, legitimación y tramitación}

El \art{169 A} precisa tres cuestiones que la práctica suele dar por supuestas:

\begin{description}[leftmargin=0pt,style=nextline,font=\sffamily\bfseries]
  \item[Oportunidad] La solicitud procede \emph{en cualquier etapa del procedimiento},
  mientras no se encuentre firme o ejecutoriada la resolución de término.

  \item[Legitimación asimétrica] El liquidador \destacar{deberá} solicitar la
  declaración de mala fe cuando concurra alguna de las causales: es un deber, no una
  facultad. Cualquier acreedor \destacar{podrá} formular la misma solicitud.

  \item[Tramitación] La solicitud se tramita \destacar{en cuaderno separado, como
  incidente}. La prueba se valora conforme a las reglas de la sana crítica y, en lo demás,
  rigen las normas del Título IX del Libro Primero del \cpc{}.
\end{description}

\begin{foco}[La consecuencia de que sea un deber del liquidador]
Si el liquidador advierte una causal y no solicita la declaración, incurre en
incumplimiento de un deber legal, con las consecuencias que ello acarrea en el control de
su cuenta y en sede disciplinaria. El acreedor que detecta la causal antes que el
liquidador tiene, por tanto, dos vías: promover él mismo el incidente, o requerir al
liquidador y dejar constancia de su omisión.
\end{foco}

\subsection{Legitimación y deber del liquidador}

Deber del liquidador de solicitar la declaración de mala fe ante la concurrencia de
indicios, y facultad de los acreedores de incoar directamente la demanda incidental.

\subsection{Valoración de la prueba y régimen recursivo}

Estándar de la sana crítica y apelación concedida en el solo efecto devolutivo.

%==============================================================================
\bloquePractico
%==============================================================================

\subsection{Modelo de demanda incidental de declaración de mala fe}

\begin{escrito}[Modelo — Incidente de declaración de mala fe del deudor]
\small
\noindent\textsc{S.J.L. en lo Civil} --- Tribunal del procedimiento de liquidación
concursal. Causa Rol [\textsc{rol}]

\medskip
\noindent [\textsc{nombre del acreedor demandante}], por sí y en representación de
[\textsc{razón social}], en los autos caratulados «[\textsc{carátula}]», a US.\
respetuosamente digo:

\medskip
Que, en conformidad a lo dispuesto en el \art{169 A}, incorporado por la \Lreforma{},
vengo en promover \textbf{incidente de declaración de mala fe del deudor}, a fin de que
se le deniegue el beneficio de extinción de los saldos insolutos de sus deudas, con
arreglo a las siguientes consideraciones de hecho y de derecho:

\begin{enumerate}[leftmargin=1.6em]
  \item \textbf{Configuración de la causal.} Conforme al numeral 1 del \art{169 A}, los
  antecedentes patrimoniales declarados por el deudor en su presentación voluntaria son
  falsos y manifiestamente incompletos.

  \item \textbf{Hechos indiciarios.} El deudor declaró un pasivo de [\textsc{monto}] y la
  posesión de un solo bien mueble valorado en [\textsc{monto}]. No obstante, de los
  antecedentes obtenidos de la \cmf{} y de la carpeta tributaria del \sii{}, que
  acompaño en el otrosí, consta que el deudor es titular de [\textsc{instrumentos de
  ahorro previsional voluntario}] por un valor total de [\textsc{monto}], y que enajenó
  un inmueble a favor de [\textsc{persona relacionada}] con anterioridad inmediata a la
  presentación concursal, omitiendo toda declaración de estas circunstancias.

  \item \textbf{Valoración bajo sana crítica.} Estos indicios objetivos, analizados a la
  luz de las máximas de la experiencia y de las reglas de la lógica, demuestran una
  conducta elusiva orientada a vaciar el patrimonio realizable en perjuicio de la masa.
\end{enumerate}

\medskip
\noindent\textbf{Por tanto}, ruego a US.\ tener por interpuesto el incidente, conferir
traslado al deudor, recibirlo a prueba bajo las reglas de la sana crítica y, en
definitiva, declarar que el deudor es de mala fe, denegándole el beneficio del
\discharge{}.
\end{escrito}

\begin{alerta}[Carga de la alegación]
El incidente no prospera con la mera afirmación de la omisión: exige acompañar la fuente
documental que acredita el activo o el acto omitido. Alegar sin prueba documental expone
al acreedor a la condena en costas.
\end{alerta}

%==============================================================================
\bloqueJurisprudencial
%==============================================================================

\subsection{Naturaleza sancionatoria civil y el debate doctrinal}

La figura, introducida por la \Lreforma{} \citeyearpar{ley-21563} para frenar la proliferación
de liquidaciones voluntarias carentes de activos, ha suscitado un debate doctrinal intenso
sobre su alcance.

\textcite{alarcon-mala-fe} ha formulado severas críticas a los efectos del inciso cuarto del
\art{169 A}, advirtiendo sobre la \destacar{incertidumbre jurídica} que introduce al carecer de
parámetros objetivos para delimitar la noción de «abuso del deudor persona natural». Por su
parte, \textcite{goldenberg-sobreendeudamiento} ha desarrollado la necesidad de examinar la
gravedad de los hechos desde la perspectiva de la \destacar{proporcionalidad y la transparencia
patrimonial}, evitando que un filtro pensado contra el fraude termine castigando la simple
desprolijidad del deudor de buena fe. Sobre la naturaleza y los límites de la figura, véanse
además \textcite{doc-mala-fe-uv} y \textcite{doc-limites-mala-fe}; sobre su proyección penal,
\textcite{doc-mala-fe-penal}.

\verificar{Sistematizar las primeras resoluciones de los juzgados civiles dictadas tras la
entrada en vigencia de la Ley 21.563.}

\subsection{¿Basta la omisión no fraudulenta?}

El punto de mayor incertidumbre práctica es si la sola omisión de un activo menor ---por
ignorancia o falta de asesoría técnica del deudor \mipe{}--- basta para denegar objetivamente
la descarga, o si la sanción exige acreditar una intención positiva de defraudar. La tensión es
la que anticipa el debate doctrinal: un estándar puramente objetivo da certeza pero castiga al
descuidado; un estándar que exija dolo protege al deudor honesto pero dificulta la prueba.

\begin{foco}[La proporcionalidad como criterio rector]
Mientras la jurisprudencia asienta criterio, la lectura que propone la doctrina
---examinar la gravedad y la proporcionalidad, no la mera existencia de una omisión--- es la
más defendible para el deudor. Una cuenta de ahorro olvidada de escaso monto no es equiparable
a la enajenación simulada de un inmueble al cónyuge: el \art{169 A} debería reservarse para lo
segundo.
\end{foco}
